% !TEX program = pdflatex
% Jaume-Vis: what a small VLM can and cannot do, and why
\documentclass[10pt,twocolumn,letterpaper]{article}

\usepackage[margin=1in]{geometry}
\usepackage{amsmath,amssymb}
\usepackage{graphicx}
\usepackage{booktabs}
\usepackage{xcolor}
\usepackage[colorlinks=true,linkcolor=blue,citecolor=blue,urlcolor=blue]{hyperref}
\usepackage{caption}
\usepackage{enumitem}
\usepackage{times}

\captionsetup{font=small,labelfont=bf}

\title{\vspace{-1.2em}What a 350M Vision-Language Model Can and Cannot Do:\\Three Negative Results and One Diagnostic Protocol}
\author{Kru Zeppesey\\\small Jaume Project\\\small ModelScope: \texttt{LightNeko/Jaume}}
\date{}

\begin{document}
\maketitle

\begin{abstract}
We report three controlled negative results from building a 350M
vision-language model, and one diagnostic protocol that explains all of them.
(i) Looped transformers do not help: across eight task structures and four
modalities, eight layers always suffice, so extra rounds repeat work rather
than adding reasoning depth. (ii) A semantic vision tower cannot read:
SigLIP2 and InternViT-448px both score 0.21 on a four-way text-location task,
at chance, while an OCR-trained encoder (TrOCR) reaches 0.75 on the same
data. (iii) Exact-match scoring hides real capability: the same model that
scores 0.11 on point-localisation scores 0.50 under a distance tolerance, a
3--5x gap. The unifying protocol is a \emph{pre-flight feasibility check}:
before spending training compute, probe the frozen encoder with a linear
classifier; if it is at chance, the information is not in the representation
and no amount of training will recover it. Applying this check turned a
projected multi-day OCR effort into a three-minute decision, and it predicts
the looping results before any loop is trained.
\end{abstract}

\section{Introduction}

Negative results are usually hard to publish because it is hard to prove that
a failure is structural rather than a tuning accident. We collected three
such results while building Jaume-Vis, a 350M parameter vision-language model
on a single Ascend 910B3, and in each case we also found a measurement that
turns the failure into a prediction.

The thread connecting them is simple. A representation either contains the
information a task needs or it does not. If it does, training will find it
given enough data; if it does not, no training helps. This sounds obvious, but
it is routinely violated: we spent two rounds of tuning on an OCR task whose
input representation had already destroyed the glyphs.

\paragraph{Contributions.}
\begin{enumerate}[leftmargin=1.2em,itemsep=1pt,topsep=2pt]
\item A depth-saturation criterion that predicts when looping cannot help,
      validated on eight task structures (Section~\ref{sec:loop}).
\item Direct evidence that semantic vision towers cannot read glyphs, and
      that resolution is not the limiting factor (Section~\ref{sec:ocr}).
\item A demonstration that exact-match metrics can under-report capability by
      3--5x, and the metric design that fixes it (Section~\ref{sec:metric}).
\item A pre-flight protocol: linear-probe the frozen encoder before training
      (Section~\ref{sec:protocol}).
\end{enumerate}

\section{Setup}

All experiments run on one Ascend 910B3 with 61\,GiB HBM. The language
backbone is Jaume, a 350M MoE transformer (16 layers, $d=512$, 4 experts,
top-2). Vision towers are frozen unless stated. Training uses AdamW with
bf16 autocast and cosine decay.

\section{Depth saturation predicts that looping cannot help}
\label{sec:loop}

\subsection{The criterion}

A looped transformer applies the same block repeatedly. If rounds are extra
depth, a loop can only help when depth helps. We therefore measure that
directly: train plain 8L and 16L models on the target task with identical data
and steps.

\begin{quote}
\begin{itemize}[leftmargin=1em,itemsep=1pt,topsep=2pt]
\item 16L beats 8L by $>0.01$: the task is depth-limited, a loop has room.
\item 16L $\approx$ 8L: the task is computed in one pass; looping cannot win.
\end{itemize}
\end{quote}

The scan costs about 0.3 GPU-hours versus 10+ for a full loop comparison.

\subsection{Results}

\begin{table*}[t]
\centering
\small
\begin{tabular}{llcccl}
\toprule
Task & Intended serial structure & 8L & 16L & Learnable & Depth-limited \\
\midrule
Chain lookup         & $k$-hop composition        & 0.0000 & 0.0000 & yes & no (closure in one pass) \\
Register program     & 16 sequential mod-13 ops   & 0.0000 & 0.0000 & yes & no (parallel width) \\
Sequence reduction   & 256-number sum/max         & 0.0000 & 0.0000 & yes & no (parallel reduce) \\
Two-stage derivation & stage 2 uses stage 1       & 0.1815 & 0.1812 & yes & no \\
Colour counting      & $3\times3$ attribute pairs & 0.983  & 0.986  & yes & no \\
Constraint path      & propagation depth $K$      & 0.0000 & 0.0000 & yes & no (transitive closure) \\
Parity / mod-count   & running mod-2/3 state      & 0.50   & 0.56   & no  & n/a (not learnable) \\
8-digit addition     & carry chain, exact match   & 0.8900 & 0.8967 & yes & no \\
\bottomrule
\end{tabular}
\caption{Depth scan across eight task structures. The last row is the
strictest: genuinely learnable, genuinely serial, strict scoring. 16L gains
only 0.0067 while costing $1.8\times$ the time.}
\label{tab:loop}
\end{table*}

Every task saturates at eight layers. 8-digit addition is the decisive case:
it is the task the loop literature uses to argue for depth, it is learnable
(0.89 exact match), it is strictly scored, and yet 8 layers match 16.

\subsection{Mechanism}

Inter-round relative distance decays geometrically:
\[
T{=}8:\; 0.552 \to 0.363 \to 0.213 \to 0.137 \to 0.097 \to 0.074 \to 0.060
\]
ratio $\approx 0.6$, and principal-subspace overlap saturates at 0.99. The
loop is a contraction to a fixed point, not staged reasoning.

\section{A semantic vision tower cannot read}
\label{sec:ocr}

\subsection{The task}

We render forms with four labelled boxes and ask which box holds a target
word. A linear probe on raw pixels reaches 100\%, so the render is separable.

\subsection{Results}

\begin{table}[t]
\centering
\small
\begin{tabular}{lccc}
\toprule
Vision tower & Params & Resolution & Accuracy \\
\midrule
SigLIP2-base   & 93M  & 224 & 0.210 \\
InternViT-300M & 304M & 448 & 0.209 \\
\textbf{TrOCR} & 334M & 384 & \textbf{0.750} \\
\bottomrule
\end{tabular}
\caption{Four-way text-location, chance 0.25. Both semantic towers are at
chance; only the OCR-trained encoder reads.}
\label{tab:ocr}
\end{table}

\subsection{It is not resolution}

InternViT runs at 448px and still scores 0.21, so the failure is not the
number of pixels. Tiling (2x2 and 3x3, each tile resized back to 224) moves
the probe from 0.179 to 0.254, i.e. to chance. The limiting factor is the
training objective, not the input bandwidth: contrastive image-text
pretraining teaches global semantic matching, never glyph discrimination.

\subsection{The fix is a different encoder, not more compute}

TrOCR's encoder, trained on image-to-text pairs, reaches 0.75 zero-shot. With
it as a second tower and a frozen Jaume backbone, a 350M model reaches 0.99 on
text location and 0.89 on character reading after training.

\section{Exact match under-reports capability}
\label{sec:metric}

A point-localisation head that must reproduce an exact digit string scores
0.109. The same predictions evaluated with a $\pm1$ grid-cell tolerance score
0.504. For boxes, exact match gives 0.058 while IoU$>0.3$ gives 0.158, and a
well-trained model reaches 0.347.

The gap is not cosmetic: it changed our conclusion from ``localisation is
broken'' to ``localisation works, the metric does not''. Detection has used
IoU for decades for exactly this reason.

\paragraph{A second metric pitfall.} Normalised text answers must be
case-folded and whitespace-collapsed before comparison. TrOCR reads
``Beauty Contest'' as ``BEAUTY CONTEST''; scoring that as wrong measures
formatting, not reading.

\section{The pre-flight protocol}
\label{sec:protocol}

We propose a three-minute check before any training run:

\begin{enumerate}[leftmargin=1.2em,itemsep=2pt,topsep=2pt]
\item Freeze the candidate encoder.
\item Extract features for a few hundred samples.
\item Fit a linear classifier.
\item If accuracy is at chance, the representation lacks the information.
      Change the encoder; do not train harder.
\end{enumerate}

For depth, the analogous check is the 8L/16L scan. Both are cheap, both are
predictive, and in our project each one replaced hours of failed training.

\section{What the model can do}

The resulting 350M checkpoint holds ten capabilities simultaneously:

\begin{table}[t]
\centering
\small
\begin{tabular}{lc}
\toprule
Capability & Score \\
\midrule
Semantic understanding (synthetic) & 0.956 \\
Real image classification (CIFAR-10) & 0.954 \\
OCR localisation (synthetic) & 0.970 \\
OCR reading (synthetic) & 0.791 \\
Document field (synthetic) & 0.612 \\
Real document reading (FUNSD) & 0.281 \\
Real document field type (FUNSD) & 0.493 \\
Object counting & 0.910 \\
Object detection (IoU$>0.3$) & 0.347 \\
Object pointing ($\pm1$ cell) & 0.645 \\
\bottomrule
\end{tabular}
\caption{Ten capabilities in one 350M checkpoint. Multi-task interference is
handled by question-conditioned gating with residual streams.}
\label{tab:matrix}
\end{table}

Two things make this possible. First, task-conditioned modulation: a gate
derived from the question embedding mixes the two visual streams, and both raw
streams stay in the sequence as residuals. Without it, OCR collapses from
0.89 to 0.65 when semantics are added. Second, explicit coordinate tokens:
counting and pointing are decoded as grid indices rather than inferred from a
feature vector.

\section{Limitations}

One model size, one hardware target, largely synthetic training data. The
depth criterion is a negative screen: saturation rules looping out, but the
absence of saturation does not guarantee a loop will help. Real-document
reading remains at 0.281 and needs a dedicated OCR model in any product.

\section{Conclusion}

Three failures, one protocol. Representations either contain the information
or they do not; a linear probe answers that question in minutes, and it
predicted every one of our results before the expensive run.

\bibliographystyle{plain}
\bibliography{refs}

\end{document}